% Lost in Pooling -- CSCI 566 midterm report
% Build: ./build.sh   (tectonic midterm.tex, run from report/)
% Figures are expected in report/figures/; every \includegraphics is guarded so the
% document compiles with a grey placeholder when a figure is missing.
% Table 1 is \input from ../results/ and is also guarded.
\documentclass[10pt,letterpaper,twocolumn]{article}

% ---------- page and fonts ----------
\usepackage[margin=0.75in]{geometry}
\usepackage[T1]{fontenc}
\usepackage{newtxtext,newtxmath}
\usepackage{microtype}
\setlength{\columnsep}{0.25in}

% ---------- packages ----------
\usepackage{amsmath}
\usepackage{graphicx}
\usepackage{booktabs}
\usepackage{xcolor}
\usepackage{caption}
\captionsetup{font=small,labelfont=bf,skip=4pt}
\usepackage{titlesec}
\titleformat{\section}{\large\bfseries}{\thesection}{0.6em}{}
\titleformat{\subsection}{\normalsize\bfseries}{\thesubsection}{0.6em}{}
\titlespacing*{\section}{0pt}{2.2ex plus 0.6ex minus 0.2ex}{1.0ex}
\titlespacing*{\subsection}{0pt}{1.6ex plus 0.4ex minus 0.2ex}{0.6ex}
\titlespacing*{\paragraph}{0pt}{1.1ex plus 0.3ex minus 0.1ex}{0.8em}
\usepackage{enumitem}
\setlist{nosep,leftmargin=*}
\usepackage[numbers,sort&compress]{natbib}
\renewcommand{\bibfont}{\footnotesize}
\setlength{\bibsep}{1pt plus 0.3ex}
\PassOptionsToPackage{hyphens}{url}
\usepackage[colorlinks=true,allcolors=black,pdftitle={Lost in Pooling: Midterm Report},pdfauthor={Sarmukaddam, Shah, Desai, Thakur, Parasher, Kale}]{hyperref}

% ---------- helper macros ----------
% Red bold placeholder for every number / detail still to be filled.
\newcommand{\todotext}[1]{{\color{red}\bfseries [TODO: #1]}}
\newcommand{\todo}[1]{\ifmmode\text{\todotext{#1}}\else\todotext{#1}\fi}
% File / identifier names with underscores, safe in captions and in \todo.
\newcommand{\fname}[1]{\texttt{\detokenize{#1}}}
% Grey placeholder box drawn when a figure file is absent.
\newcommand{\figplaceholder}[2]{%
  \colorbox{black!12}{\parbox[c][1.9in][c]{\dimexpr#2-2\fboxsep\relax}{%
    \centering\small\color{black!60} figure not found\\[3pt]
    {\footnotesize\ttfamily\expandafter\detokenize\expandafter{#1}}}}}
% \guardedfig{path}{width}: includes the figure if it exists, otherwise the placeholder.
\newcommand{\guardedfig}[2]{\IfFileExists{#1}{\includegraphics[width=#2]{#1}}{\figplaceholder{#1}{#2}}}

% Paths to figure and table hooks (edit here if names change)
\newcommand{\figBlindness}{figures/blindness_curve__dinov2_vitb14_reg__res224.pdf}
\newcommand{\figAttention}{figures/attention_examples.png}
\newcommand{\tabBlindness}{../results/blindness_table__dinov2_vitb14_reg__res224.tex}

\raggedbottom
\begin{document}

% ======================================================================
%  Title block and abstract, single column
% ======================================================================
\makeatletter
\twocolumn[
\begin{@twocolumnfalse}
\begin{center}
  {\LARGE\bfseries Lost in Pooling: Recovering Small-Signal Evidence\\[3pt]
   Discarded by Frozen Foundation Models\par}
  \vspace{5pt}
  {\large\bfseries Midterm Report\par}
  \vspace{12pt}
  {\normalsize Dhairya Shah \quad Raghav Sarmukaddam \quad Yash Desai \quad
  Prisha Thakur \quad Ansh Parasher \quad Advait Kale\par}
  \vspace{4pt}
  {\small\ttfamily \{dhairyad, sarmukad, ysdesai, prishath, aparashe, advaitka\}@usc.edu\par}
  \vspace{4pt}
  {\small University of Southern California \ \textbullet\ \ CSCI 566, Fall 2026\par}
  \vspace{16pt}
  {\normalsize\bfseries Abstract\par}
  \vspace{4pt}
\end{center}
\begin{list}{}{\setlength{\leftmargin}{0.55in}\setlength{\rightmargin}{0.55in}}\item\relax
\noindent
The dominant transfer recipe freezes a large pretrained backbone, keeps its single global
embedding, and trains a light classifier on top. Recent evidence from chest radiography shows
that this recipe fails silently on small targets: global (CLS) embeddings of five foundation
models score at chance on small lesions while the same models' patch tokens retain near-perfect
signal once a region of interest is supplied. The evidence survives the encoder and is lost at
the aggregation step. Our pre-proposal asked how general this is, what causes it, and whether
it can be recovered without extra supervision. Following the instructor's feedback we narrowed
the midterm to one controlled comparison: one natural-image dataset (MS-COCO 2017, 80-way
multi-label) and one frozen backbone (DINOv2 ViT-B/14 with registers), with every read-out
trained and evaluated on identical cached tokens. We stratify per-class AUROC by the relative
area of the largest instance in six log-spaced bins, two of which lie below the area of a single
patch, and compare CLS, mean, concatenated, max and max-MIL read-outs against a
told-where-to-look oracle that pools only the patches inside ground-truth boxes. We also train
a gated attention MIL head from image-level labels alone whose entropy-regularised limit is
exactly mean pooling, so the baseline is nested in the model, and we audit whether its attention
lands on the object. On 5,000 COCO test images the best global read-out falls from 0.99 AUROC
on large targets to 0.86 on sub-patch targets, while the oracle on the same tokens stays at
0.97, so the evidence survives the encoder and is lost at pooling; the head recovers a third of
that gap (0.90) and lifts mAP from 0.73 to 0.80, with attention mass on the target 12 to 24
times chance. We explain
why natural images, with their strong context priors, may show a shallower cliff than
radiographs, and set out the three-week plan to the final report.
\end{list}
\vspace{14pt}
\end{@twocolumnfalse}
]
\makeatother

% ======================================================================
\section{Introduction}
% ======================================================================
Transfer learning with a frozen foundation model is cheap, reproducible and usually good
enough: compute an embedding once, fit a linear probe, ship. The hidden assumption is that the
embedding keeps whatever the downstream task needs. A recent evaluation of five chest-radiograph
foundation models shows that this assumption can fail completely for small targets. Linear
probes on the global (CLS) embeddings score at chance on small lesions, AUC 0.500 to 0.524,
while probes on the same models' patch tokens restricted to a supplied region of interest are
near perfect \citep{smalllesion2026}. The signal survives the encoder and is destroyed when the
tokens are collapsed into one vector.

Our pre-proposal asked three questions: how general is this effect, what causes it, and can it
be recovered without extra supervision? The feedback we received (Section~\ref{sec:feedback})
was that the plan, with seven backbones, several modalities, mechanism ablations and
distillation, was too large for one semester, and that the causal claim needed a controlled
experiment rather than similar curves across modalities. This report describes the narrowed
project. We take one natural-image dataset, MS-COCO 2017 \citep{lin2014coco}, and one frozen
backbone, DINOv2 ViT-B/14 with registers \citep{oquab2024dinov2,darcet2024registers}, cache
the CLS and patch tokens once, and train every read-out on the same cached tokens. We stratify
per-class AUROC by the relative area of the largest instance of the class in six log-spaced
bins and compare the global read-outs against an oracle that is told where to look. We then
ask whether a small gated attention head trained from image-level labels alone can close the
gap, and whether its attention actually lands on the object.

Section~\ref{sec:lit} reviews the related literature, Section~\ref{sec:progress} describes
progress since the pre-proposal and our response to feedback, Section~\ref{sec:setup} fixes
the experimental protocol, Section~\ref{sec:results} reports initial results, and
Section~\ref{sec:plan} gives the plan to the final report. Code and results are at
\url{https://github.com/DhairyaShah5/lost-in-pooling}.

% ======================================================================
\section{Literature Review}
\label{sec:lit}
% ======================================================================

\subsection{Frozen features and linear probing}
Freezing a pretrained encoder and training a linear classifier on its output is the oldest and
still the most common way to use a foundation model. It is the standard evaluation for
self-supervised vision models \citep{caron2021dino,oquab2024dinov2}, it is how CLIP reports
transfer on its suite of downstream datasets \citep{radford2021clip}, and
\citet{kornblith2019transfer} showed that ImageNet accuracy predicts linear-probe transfer
across a wide range of target tasks. The recipe is attractive because features can be cached
once and probes are cheap, so a hospital or a lab can evaluate a backbone on a laptop. The
implicit assumption is that the single vector handed to the probe retains whatever the task
needs. For whole-image categories this holds well: a frozen DINOv2 ViT-g/14 with a linear
layer exceeds 86\% top-1 on ImageNet \citep{oquab2024dinov2}. Whether it holds when the label
depends on a few percent of the pixels is the question this project asks.

\subsection{How a ViT produces one vector}
A vision transformer emits one token per patch plus, in the original design, a learned class
token whose final state is read out by the classification head \citep{dosovitskiy2021vit}.
Three global read-outs are in common use. The CLS token aggregates through self-attention
inside the network, so what it keeps is shaped entirely by the pretraining objective. Global
average pooling (GAP) of the patch tokens is the alternative; \citet{beyer2022betterplain}
found that a plain ViT trained with GAP matches or slightly exceeds the CLS variant on
ImageNet, and \citet{raghu2021vitcnn} showed that CLS- and GAP-trained models organise spatial
information differently in their upper layers. The third option is a learned multihead
attention pooling (MAP) head, a single cross-attention layer with a learned query, introduced
for scaling ViTs \citep{zhai2022scaling} and used by SigLIP \citep{zhai2023siglip}. All three
produce one fixed-size vector before any downstream task is specified, so none of them can know
which patches will matter for a later label. \citet{darcet2024registers} add a complication:
large supervised and self-supervised ViTs recycle low-information background patches as
high-norm tokens that carry global rather than local content, which corrupts the patch map;
adding dedicated register tokens removes the artefacts and improves dense downstream tasks. We
therefore use the register variant of DINOv2, since the patch tokens are our object of study.

\subsection{Patch tokens localise objects without supervision}
Self-supervised ViT features are unusually good at separating foreground from background with
no labels. \citet{caron2021dino} observed that the CLS attention maps of DINO segment the main
object. LOST \citep{simeoni2021lost} localises an object in a single forward pass by picking
the patch with the fewest positive correlations to the others as a seed and growing a region of
correlated patches, with no training at all. TokenCut \citep{wang2022tokencut} builds a graph
over patch tokens and applies a normalised cut to separate the object from the background.
\citet{amir2022deepvit} show that intermediate DINO-ViT features act as dense descriptors with
part-level correspondence across images. DINOv2 extended this with large-scale curated
pretraining and reports strong linear segmentation and depth results from frozen patch tokens
\citep{oquab2024dinov2}. Together these results say that the information needed to find and to
classify an object survives in the tokens; what is at issue is whether global pooling keeps it.

\subsection{Small objects, scale and resolution}
Small objects are the hardest case for detectors. \citet{singh2018snip} showed that scale
variation dominates detection error and proposed scale normalisation, training only on objects
near the pretraining scale. Slicing inference \citep{akyon2022sahi} recovers small objects by
running the detector on overlapping tiles, which raises the effective resolution at the cost of
several forward passes. \citet{cheng2023smallobject} survey small-object detection and release
benchmarks stratified by absolute object size. For ViTs the relevant knob is the patch size:
FlexiViT \citep{beyer2023flexivit} trains one model across patch sizes and shows a smooth
trade-off between compute and accuracy, and DINOv2 raises its input resolution to 518 px for the
final epochs to sharpen patch features \citep{oquab2024dinov2}. All of these methods act before
or inside the encoder. Our question is complementary: given a fixed encoder at a fixed
resolution, how much of the small-object signal is lost \emph{after} the encoder, at the
aggregation step? We treat input resolution as a planned ablation rather than the main variable.

\subsection{Multiple-instance learning pooling}
Multiple-instance learning (MIL) treats an image as a bag of instances with a single bag label
\citep{dietterich1997mil}. The pooling function is the main design choice. Max pooling selects
one instance and encodes the classic MIL assumption that a positive bag contains at least one
positive instance; mean pooling spreads responsibility over the bag;
\citet{wang2018revisiting} compare instance-level and embedding-level MIL networks and find
embedding-level pooling more robust. \citet{ilse2018attention} proposed attention-based MIL,
where a two-layer network with a gated $\tanh$/sigmoid nonlinearity produces instance weights
that sum to one and the bag embedding is the weighted sum, so the weights double as an
interpretable map. This is now the standard aggregator over frozen tile features in
computational pathology: \citet{campanella2019clinical} train a clinical-grade classifier with
max-MIL over tile features, and CLAM \citep{lu2021clam} uses class-specific gated attention
over frozen tile embeddings with instance-level clustering as a regulariser. A whole-slide image
is the extreme case of our problem, a tumour of a few tiles in a bag of tens of thousands, and
the pathology literature already contains the pattern we want to test on natural images: frozen
encoder, learned pooling, image-level labels only.

\subsection{Size-stratified evaluation and the motivating paper}
COCO reports average precision separately for small (area below $32^2$ pixels), medium and
large objects \citep{lin2014coco}, and most detection papers follow;
\citet{cheng2023smallobject} argue for finer stratification. Classification benchmarks almost
never stratify by target size. In medical imaging this matters because the clinically important
finding is often small. ChestX-ray8 \citep{wang2017chestxray8} introduced weakly supervised
classification and localisation of thorax diseases from image-level labels, and foundation
models for the domain now include BiomedCLIP \citep{zhang2023biomedclip} and RAD-DINO
\citep{perezgarcia2025raddino}. The paper that motivates us evaluated frozen global embeddings
of five such models on small lesions and found AUC at chance (0.500 to 0.524), while patch
tokens inside a supplied region of interest retained near-perfect signal
\citep{smalllesion2026}; it concludes that pre-deployment evaluation must stratify by lesion
size. Whether the effect is specific to radiographs, where the lesion is subtle and the rest of
the image is uninformative about it, or a general property of global pooling, is open. Context
is the obvious confound on natural images: objects co-occur with predictable scenes and
companions, and recognition systems, human and artificial, exploit this
\citep{oliva2007context}.

\subsection{How our work differs}
Each of the lines above answers part of our question. What is missing is a controlled comparison
that holds the encoder, the tokens and the training images fixed, varies only the read-out, and
reports performance as a function of target size. We provide that comparison on one dataset and
one backbone, with a told-where-to-look oracle on the same tokens as the ceiling. We use an
attention head that contains mean pooling as a special case, so the baseline is nested in the
model and any gain is attributable to where the head looks. Finally, we audit the head's
attention against ground-truth boxes, which separates ``sees the object'' from ``guesses from
context'' in a way that the AUROC curve alone cannot. Audio event detection is the natural second modality, where a ``small'' event is a brief one
and attention pooling over frames is already standard \citep{kong2019attention,chen2023beats};
per the feedback it is deferred until the vision result is settled.

% ======================================================================
\section{Progress So Far}
\label{sec:progress}
% ======================================================================

\subsection{Response to pre-proposal feedback}
\label{sec:feedback}
The pre-proposal was graded 5/5 with four pieces of feedback. We list each and what we changed.

\begin{itemize}
\item \textbf{Scope.} Seven backbones, several modalities, mechanism ablations and distillation
  were judged too much for one semester, with the advice to do the pooling comparison first on
  one image dataset and one backbone using the same frozen tokens, then expand. We did exactly
  this. The midterm study is MS-COCO 2017 with DINOv2 ViT-B/14 (registers) at 224 px. Audio,
  distillation and additional backbones are deferred. The chest-radiograph replication is
  scheduled for the final three weeks (Section~\ref{sec:plan}) and only if the COCO result is
  in. LUNA16 was dropped outright: it is 3D CT, the notion of relative box area does not carry
  over cleanly to volumes, and it would not answer the question any better than COCO.
\item \textbf{The causal claim.} Similar curves across modalities do not show that pooling
  caused the failure. Our design is now a within-subjects comparison: identical cached tokens,
  identical train/validation/test images, identical optimiser and early-stopping rule; only the
  read-out changes. Any difference in the size curve is therefore attributable to the read-out,
  and the oracle read-out on the same tokens bounds how much signal the tokens carry at each
  size. We will not use cross-modality similarity as evidence of mechanism.
\item \textbf{The size-based evaluation is useful on its own.} We agree and built the
  evaluation harness as a standalone module with a file contract, so that the blindness curve
  for the fixed read-outs is a deliverable independent of whether the learned head helps.
\item \textbf{Keep negative findings specific.} Every result in this report is stated for
  DINOv2 ViT-B/14 with registers, at 224 px input, on COCO val2017, with probes trained on a
  9,000-image subset of train2017. We do not generalise beyond that until the replication runs.
\end{itemize}

\subsection{Design decisions and their reasons}
\paragraph{Dataset: MS-COCO 2017.}
Every non-crowd instance has a bounding box, so relative size is defined for every positive;
80 classes give enough per-bin positives for a macro average; roughly 40\% of instances are
``small'' under the official definition \citep{lin2014coco}; and the data are public, which
keeps the pipeline reproducible by anyone with a laptop. The known cost is context priors
(Section~\ref{sec:discussion}), which we address with the oracle gap and the localisation
audit rather than by pretending they do not exist.

\paragraph{Backbone: DINOv2 ViT-B/14 with registers, frozen.}
Self-supervised, so no COCO label leaks into the encoder; patch tokens that are known to
localise objects; register tokens that remove the high-norm artefacts
\citep{darcet2024registers}; and a ViT-B that fits the compute we have, a single
Apple-silicon laptop with 16 GB of unified memory (measured throughput: 38 images per second at 224 px).

\paragraph{Input handling: squash-resize to $224\times224$, no crop.}
A centre crop would discard boxes near the border; squashing keeps every box inside the token
grid at the price of a distortion that affects every read-out equally. The backbone runs once
per image and every read-out, including the attention head, reads the same fp16 cache, which
is what makes the comparison same-token by construction.

\subsection{Pipeline and repository}
The repository is organised as a small library \fname{src/lip/} with numbered scripts under \fname{scripts/} that run in order, a
written data and feature contract (\fname{docs/CONTRACTS.md}) that every module reads and
writes against, and committed JSON results under \fname{results/}. The steps are:
\begin{enumerate}
\item \textbf{Download.} COCO 2017 annotations, all 5,000 val2017 images, and a seeded
  10,000-image subset of train2017 (sorted image ids, permuted with seed 0, first 10,000).
\item \textbf{Manifest, labels and boxes.} One manifest row per image with split and size;
  an $N\times80$ presence matrix, an $N\times80$ matrix of largest-instance box-area fractions
  and the same for segmentation area; and a boxes table in normalised coordinates. Crowd
  annotations are ignored throughout.
\item \textbf{Feature cache.} DINOv2 ViT-B/14-reg on the MPS backend, squash-resize and
  ImageNet normalisation, fp16 memmap writer, one directory per backbone, resolution and split.
\item \textbf{Fixed read-outs.} CLS, mean, CLS+mean, per-dimension max, max-MIL and oracle
  region pooling, each with a linear probe trained by joint binary cross-entropy with early
  stopping on the 1,000-image validation split.
\item \textbf{Learned head.} Gated attention MIL head over the memmapped tokens with an
  entropy regulariser, plus the localisation metrics.
\item \textbf{Evaluation and figures.} Per-class per-bin AUROC, macro average, bootstrap
  confidence intervals, overall AUROC and mAP, results JSON, the blindness-curve figure and the
  per-bin table.
\end{enumerate}
Every script is seeded, takes \texttt{--backbone}, \texttt{--res} and \texttt{--out} flags, and
writes to the paths in the contract. Unit tests run on synthetic data without the real dataset. The code is public at
\url{https://github.com/DhairyaShah5/lost-in-pooling}.
Status at submission: steps 0 to 5 are complete and committed: the data (15,000 images), both
feature caches (224 px for all splits; 448 px for test, val and 5,000 training images), the six
fixed read-outs, the oracle and its negative-sampling ablation, the attention head with a
four-value $\lambda$ sweep, the localisation audit, the resolution ablation, and every figure
and table in this report, each regenerated from the committed \texttt{results/*.json} files by
\texttt{scripts/07\_plots.py}. Not yet run: the multi-scale head, top-$k$ max-MIL variants, and
any second backbone or dataset.

% ======================================================================
\section{Experimental Setup}
\label{sec:setup}
% ======================================================================

\subsection{Data and splits}
Test set: all 5,000 COCO val2017 images. Training data: a seeded 10,000-image subset of
train2017, split 9,000/1,000 into \emph{train} (probe fitting) and \emph{val} (model selection
and early stopping only). No test image is seen during training or selection. Labels are
80-way multi-label presence; category index follows ascending COCO category id.

\paragraph{Chest radiographs (NIH ChestX-ray14).}
To test the motivating domain we add NIH ChestX-ray14 \citep{wang2017chestxray8}: 14 findings as
image-level labels, and 984 radiologist boxes for 8 findings on 880 images. We downloaded
23,839 images (all 880 boxed ones plus a random subset). Test = the 880 boxed images; train/val
= 20,721/2,238 images, patient-disjoint from test (seeded). Size = box area over the $1024^2$
frame. A pair that carries a finding label but no box counts as a positive overall and enters
no size bin. Because the training images have no boxes, the oracle cannot be fitted on them;
Section~\ref{sec:cxr-protocol} gives the cross-validated protocol.
\subsection{Backbones and feature cache}
\label{sec:backbones}
Primary: DINOv2 ViT-B/14 with four register tokens, frozen, input squash-resized to
$224\times224$ with ImageNet normalisation; we keep the final-layer-norm CLS token
$z_0\in\mathbb{R}^{768}$ and the $P=256$ patch tokens in row-major grid order, and discard the
register tokens. To test whether the pattern depends on the pretraining objective or the
patch size we add two backbones with a 16-pixel patch ($14\times14=196$ tokens at 224 px, one
patch $=0.51\%$ of the image): CLIP ViT-B/16 \citep{radford2021clip}, language-supervised,
whose global vector is the post-layer-norm CLS token and whose patch tokens we take after the
same layer norm; and SigLIP ViT-B/16 \citep{zhai2023siglip}, which has no CLS token and whose
global vector is a learned multihead attention-pooling (MAP) head, so its ``CLS'' read-out is
already a trained attention pool over the patches. For the radiographs we use DINOv2 and
RAD-DINO \citep{perezgarcia2025raddino}, a DINOv2-B/14 continued-pretrained on chest X-rays
(no registers). Every backbone is cached once per image as fp16 (global token plus patch
tokens), and all read-outs of a backbone read the same cache.
\subsection{Size variable, bins and box-to-patch mapping}
\label{sec:bins}
For class $c$ in image $n$ the size variable is $s_{n,c}=\max_{k}\,\mathrm{area}(b_k)/(W H)$
over the non-crowd instances $b_k$ of class $c$, binned at 0.1\%, 0.3\%, 1\%, 3\% and 10\%
of the image area ($<0.1\%$, 0.1--0.3\%, 0.3--1\%, 1--3\%, 3--10\%, $>10\%$). A patch is \emph{inside} a box if its centre lies in the
box; a sub-patch box that contains no patch centre is assigned the single patch containing the
box centre, so every positive has at least one inside patch.

\subsection{Read-outs compared on identical tokens}
Each read-out maps the cached tokens of an image to 80 class scores. With $\bar z =
\frac{1}{P}\sum_i z_i$:
\begin{align*}
\text{CLS:}&\quad s_c = w_c^{\top} z_0 + b_c,\\
\text{mean:}&\quad s_c = w_c^{\top} \bar z + b_c,\\
\text{CLS+mean:}&\quad s_c = w_c^{\top} [z_0;\bar z] + b_c,\\
\text{max:}&\quad s_c = w_c^{\top} \big(\max_i z_{i,\cdot}\big) + b_c \ \text{(per dimension)},\\
\text{max-MIL:}&\quad s_c = \max_i \big(w_c^{\top} z_i + b_c\big),\\
\text{oracle:}&\quad s_c = w_c^{\top}\Big(\tfrac{1}{|B_{n,c}|}\textstyle\sum_{i\in B_{n,c}} z_i\Big) + b_c .
\end{align*}
The first four are linear probes on a fixed global vector. Max-MIL scores every patch with the
same linear scorer and keeps the maximum, the classic instance-level MIL read-out. The oracle
pools only the patches $B_{n,c}$ inside the ground-truth boxes of class $c$; for negative
images, which have no box, we draw a random box whose area is sampled from that class's
positive box-area distribution, so the probe cannot learn box statistics and the oracle
measures only whether the tokens inside a region of the right size carry the class. It is the
told-where-to-look upper bound on the same tokens.

\subsection{Gated attention MIL head}
Following \citet{ilse2018attention} with class-specific attention as in CLAM
\citep{lu2021clam}, the head computes for each patch a gated hidden state
$h_i = \tanh(Vz_i)\odot\sigma(Uz_i)$, $h_i\in\mathbb{R}^{L}$, class-specific attention
$a_{i,c} = \mathrm{softmax}_i\big(u_c^{\top} h_i\big)$, a class-specific bag embedding
$g_c=\sum_i a_{i,c}\, z_i$ and a score $s_c = \theta_c^{\top} g_c + b_c$. With $L=256$ the
head has 475,728 parameters and sees image-level labels only. The loss is joint binary
cross-entropy plus an entropy term on the attention of the classes present in the image,
\begin{equation*}
\mathcal{L} = \mathrm{BCE}(s, y) + \lambda\,\frac{1}{|C_n|}\sum_{c\in C_n}\frac{H(a_{\cdot,c})}{\log P},
\end{equation*}
where $H$ is the Shannon entropy, so the term lies in $[0,1]$ and $\lambda>0$ discourages the
degenerate uniform solution. That uniform solution, $a_{i,c}=1/P$ for all $i$, gives
$g_c=\bar z$, so the head reduces exactly to the mean-pooling probe: mean pooling is a nested
special case that the head can always fall back to (it is the $\lambda\to-\infty$ limit). We
sweep $\lambda\in\{0, 0.01, 0.1\}$.

\subsection{Evaluation protocol}
For class $c$ and bin $b$, AUROC is computed with positives the test images whose largest
instance of $c$ falls in $b$ and negatives \emph{all} test images without class $c$; the
negatives are shared across bins so that bins differ only in the positives. The per-bin score
is the macro average over classes with at least 10 positives in that bin. A 95\% confidence
interval comes from a bootstrap over images (500 draws). We also report overall macro AUROC
and mAP with all sizes pooled. For the attention head we add two localisation metrics per bin:
attention mass inside the ground-truth boxes, $m_{n,c}=\sum_{i\in B_{n,c}} a_{i,c}$, compared
with its chance level $|B_{n,c}|/P$, and pointing-game accuracy, the fraction of positives
whose arg-max patch lies inside a box.

\subsection{Training details}
Linear probes are trained with AdamW (learning rate $10^{-3}$, weight decay $10^{-4}$, batch
512, 40 epochs, cosine schedule) on standardised features, keeping the epoch with the best
validation mAP; seed 0 throughout. Max-MIL uses the same optimiser with batch 128 for up to 15
epochs with patience 3. The oracle probes are per-class logistic regressions (L-BFGS, $C=1$)
with train negatives subsampled to 3,000 per class. The attention head is trained with AdamW
(learning rate $10^{-3}$, weight decay $10^{-4}$, batch 128, up to 20 epochs, patience 5) with
the same selection rule. All experiments run on one Apple-silicon laptop (M4, 16 GB). Feature
extraction runs at 38 images/s at 224 px and 10 images/s at 448 px, so the 224 px cache of all
15,000 images takes about seven minutes and occupies 5.5 GB in fp16; every probe then trains
in well under a minute. We also cache 448 px features ($32\times32$ grid, one patch
$=0.098\%$ of the image) for the test split and a 5,000-image training subset, to test whether
the sub-patch bins recover once the object spans at least one token.

\subsection{Chest radiograph protocol}
\label{sec:cxr-protocol}
The NIH training images carry labels but no boxes, so two protocols are run. (i) Image-level
training as on COCO: CLS, mean, max, max-MIL and the attention head are trained on the 20,721
training images and evaluated per size bin on the 880 boxed images. (ii) Cross-validation on
the boxed images, following the motivating study: CLS, mean, max and the oracle are each
fitted as per-class logistic regressions with 5-fold patient-disjoint cross-validation over
the 880 boxed images, so all four see identical training images and the oracle gap is a
same-data comparison; positives are boxed pairs, negatives are images without the finding
label, and unboxed positives are excluded from fitting. With 79 to 180 boxed images per
finding the bins are thin, so for this arm we lower the per-bin threshold to 5 positives per
class and report the number of contributing findings.

% ======================================================================
\section{Initial Results}
\label{sec:results}
% ======================================================================

\begin{table*}[t]
\centering\small
\caption{Macro AUROC per size bin (95\% bootstrap CI) and overall AUROC/mAP for each read-out
on identical DINOv2 ViT-B/14-reg tokens at 224 px, COCO val2017. Generated by the evaluation
script into \fname{results/blindness_table__dinov2_vitb14_reg__res224.tex}; the best
non-oracle read-out per column is in bold.}
\label{tab:blindness}
\IfFileExists{\tabBlindness}{\input{\tabBlindness}}{%
  \colorbox{black!12}{\parbox[c][0.9in][c]{0.9\textwidth}{\centering\small\color{black!60}
  table not found\\[3pt]{\footnotesize\ttfamily\detokenize{../results/blindness_table__dinov2_vitb14_reg__res224.tex}}}}}
\end{table*}

\subsection{Dataset statistics}
Table~\ref{tab:stats} gives the number of positive (image, class) pairs per size bin on the
test split and the number of classes that clear the 10-positive threshold, which determines
how many classes enter the macro average in each bin.

\begin{table}[t]
\centering\small
\caption{Positives per size bin on COCO val2017 (test). \emph{Pairs}: (image, class) pairs
whose largest instance falls in the bin. \emph{Classes}: classes with at least 10 such pairs,
i.e.\ those entering the macro average. The 9,000 training images contribute 26,202 pairs
with the same shape (1,160 in the smallest bin).}
\label{tab:stats}
\begin{tabular}{lrrr}
\toprule
Bin & Pairs & Classes & Share of pairs \\
\midrule
$<0.1\%$     & 637   & 19 & 4.4\%  \\
0.1--0.3\%   & 1,149 & 30 & 7.9\%  \\
0.3--1\%     & 2,037 & 50 & 13.9\% \\
1--3\%       & 2,324 & 62 & 15.9\% \\
3--10\%      & 3,045 & 68 & 20.8\% \\
$>10\%$      & 5,439 & 71 & 37.2\% \\
\midrule
All          & 14,631 & 80 & 100\% \\
\bottomrule
\end{tabular}
\end{table}

\subsection{The blindness curve}
Figure~\ref{fig:blindness} plots macro AUROC against size bin for every read-out on the same
tokens, and Table~\ref{tab:blindness} gives the numbers with confidence intervals.

\begin{figure}[t]
\centering
\guardedfig{\figBlindness}{\columnwidth}
\caption{Blindness curve for DINOv2 ViT-B/14-reg at 224 px on COCO val2017: macro AUROC per
size bin for each read-out trained on identical cached tokens. Shaded bands are 95\%
bootstrap intervals over images. The dotted vertical line marks the area of one patch
(0.39\%). The oracle is drawn in black; the learned head shares the axes with the fixed
read-outs. Per-bin class counts ($n$) are printed under the axis.}
\label{fig:blindness}
\end{figure}

\paragraph{Global read-outs.}
All four global read-outs fall monotonically with target size. Mean pooling, the best of them,
scores 0.992 in the $>10\%$ bin and 0.862 in the $<0.1\%$ bin, a drop of 0.13 AUROC; CLS is
lower throughout (0.978 down to 0.832, a drop of 0.15), and CLS+mean (0.990 to 0.843) and
per-dimension max pooling (0.991 to 0.846) sit between the two. The decline begins before the
one-patch boundary: in the 0.3--1\% bin, where the object spans one to two patches, every
global read-out is already about 0.06 below its large-object value. No read-out approaches
chance in any bin, unlike the chest-radiograph cliff. Contrary to what the MIL literature
would predict, max pooling does not hold up better than mean pooling on sub-patch targets
(0.846 vs.\ 0.862), and max-MIL, whose per-patch scorer should in principle detect a single
firing patch, is marginally below mean pooling there (0.851) and clearly worse on large
targets (0.945); its validation mAP peaked after the first epoch, so with a linear scorer and
9,000 training images the top-1 instance rule is a harder estimation problem than averaging,
not a better aggregator.

\paragraph{Oracle gap.}
The oracle, which pools only the patches under the ground-truth boxes, stays above 0.97 in
every bin: 0.973 in the $<0.1\%$ bin and 0.989 in the 0.1--0.3\% bin, against 0.862 and 0.918
for mean pooling. The gap between the oracle and the best global read-out is therefore
0.11 AUROC for sub-patch targets, 0.06 for targets of one to two patches, and under 0.01 for
targets above 10\% of the image. Because the oracle reads the same cached tokens as the global
probes, the evidence for sub-patch objects survives the encoder at 224 px and is lost at
aggregation, which is the controlled comparison the pre-proposal feedback asked for. The
matched-negative design matters: if negatives are instead given the global mean of all
patches (\texttt{oracle\_globalneg} in \texttt{results/}), the oracle rises to 0.988 on
sub-patch targets and 0.992 mAP overall (vs.\ 0.944), because the probe can then separate
box-pooled from globally pooled statistics without reading content. We report the matched
version as the upper bound.

\paragraph{Resolution ablation.}
If the loss were in the encoder, giving it more pixels per object should help every read-out.
We re-extracted the test split and a 5,000-image training subset at 448 px ($32\times32$ grid,
one patch $=0.098\%$, so the $<0.1\%$ bin is now about one patch and the 0.1--0.3\% bin spans
one to three patches) and trained the same probes; the 224 px probes in
Table~\ref{tab:resolution} are retrained on the same 5,000 images so that only the resolution
differs. The oracle improves in every small bin (0.973 to 0.994 for sub-patch targets), so the
tokens at 448 px carry more evidence about small objects. The global read-outs do not move:
CLS, mean and max change by at most 0.008 in the three smallest bins. Four times as many tokens
now describe the same object, and four times as many describe everything else, so the pooled
vector dilutes the object by the same factor. For DINOv2 ViT-B/14 on COCO, resolution fixes the
encoder's view of small objects but not the aggregation's.

\begin{table}[t]
\centering\small
\caption{Resolution ablation, macro AUROC on COCO val2017. Both columns per bin use probes
trained on the same 5,000 training images; the oracle uses all available training positives
(9,000 images at 224 px, 5,000 at 448 px), which if anything disadvantages 448 px.}
\label{tab:resolution}
\begin{tabular}{l cc cc cc}
\toprule
 & \multicolumn{2}{c}{$<0.1\%$} & \multicolumn{2}{c}{0.1--0.3\%} & \multicolumn{2}{c}{0.3--1\%} \\
\cmidrule(lr){2-3}\cmidrule(lr){4-5}\cmidrule(lr){6-7}
Read-out & 224 & 448 & 224 & 448 & 224 & 448 \\
\midrule
CLS    & 0.824 & 0.828 & 0.888 & 0.890 & 0.909 & 0.914 \\
Mean   & 0.853 & 0.858 & 0.917 & 0.921 & 0.928 & 0.936 \\
Max    & 0.843 & 0.837 & 0.905 & 0.906 & 0.918 & 0.927 \\
\midrule
Oracle & 0.973 & \textbf{0.994} & 0.989 & \textbf{0.996} & 0.992 & \textbf{0.999} \\
\bottomrule
\end{tabular}
\end{table}

\subsection{Learned head and localisation}
\begin{figure*}[t]
\centering
\guardedfig{\figAttention}{\textwidth}
\caption{Class-specific attention of the gated head ($\lambda=0$) for six test images whose
largest instance of the named class covers 0.15--0.30\% of the image (sub-patch to one patch).
Top: the squash-resized input with the ground-truth boxes of that class in white. Bottom: the
class's attention map on the $16\times16$ grid, normalised by its maximum. The head finds the
surfer, the traffic light, the glove and the clock; it lands one patch beside the bird; and for
the tennis racket it attends to the player's body, a context cue, not the target.}
\label{fig:attention}
\end{figure*}

\begin{table}[t]
\centering\small
\caption{Localisation audit of the attention head per size bin: mean attention mass inside the
ground-truth boxes versus its chance level (box area as a fraction of the grid), and
pointing-game accuracy (arg-max patch inside a box). $\lambda=0$; $n$ is the number of
positive (image, class) pairs. \emph{Uniform} attention would score exactly the chance column.}
\label{tab:localisation}
\begin{tabular}{lrrrrr}
\toprule
Bin & $n$ & Mass in box & Chance & Ratio & Pointing \\
\midrule
$<0.1\%$     & 637   & 0.088 & 0.007 & 12.1 & 0.129 \\
0.1--0.3\%   & 1,149 & 0.195 & 0.008 & 23.5 & 0.260 \\
0.3--1\%     & 2,037 & 0.304 & 0.013 & 23.7 & 0.382 \\
1--3\%       & 2,324 & 0.447 & 0.030 & 15.0 & 0.515 \\
3--10\%      & 3,045 & 0.636 & 0.086 & 7.4  & 0.715 \\
$>10\%$      & 5,439 & 0.813 & 0.410 & 2.0  & 0.842 \\
\midrule
All          & 14,631 & 0.567 & 0.178 & 3.2 & 0.623 \\
\bottomrule
\end{tabular}
\end{table}

\paragraph{Learned head.}
Figure~\ref{fig:blindness} and Table~\ref{tab:blindness} include the gated attention head at
$\lambda=0$. It beats every fixed read-out in every bin below 10\%: 0.899 on sub-patch targets
against 0.862 for mean pooling and 0.973 for the oracle, closing a third of the oracle gap
there, a fifth in the 0.1--0.3\% bin and a third in the 0.3--1\% bin; it raises overall mAP from
0.734 (mean) and 0.758 (CLS+mean) to 0.800. On targets above 10\% it sits 0.006 below mean
pooling, the price of a class-specific attention that no longer averages over the whole object.
Table~\ref{tab:lambda} gives the entropy sweep. A mild sharpening term ($\lambda=0.01$)
changes nothing; a strong one ($\lambda=0.1$) collapses the attention onto a few patches
(normalised entropy 0.06) and both AUROC and mAP fall. Pushing the other way
($\lambda=-0.1$) keeps the attention near uniform (entropy 0.90) and gives the best
per-bin AUROC of the sweep (0.908 on sub-patch targets) at slightly lower mAP (0.790), but its
attention mass on small targets drops to 3.6 times chance: it is mean pooling with a mild,
class-specific tilt, and it ranks images well without looking at the object much. The
$\lambda=0$ head is the one that both improves on mean pooling and localises, so we report it
as the main result; the trade-off between ranking and localisation is itself a finding we will
pursue with the multi-scale variant.

\begin{table}[t]
\centering\small
\caption{Entropy-regulariser sweep for the attention head, test split. \emph{Ent.}: normalised
attention entropy at the selected epoch (uniform $=1$). \emph{Mass}: attention mass inside
the boxes of sub-patch targets (chance 0.007). Mean pooling and the oracle are repeated for
reference.}
\label{tab:lambda}
\setlength{\tabcolsep}{3.5pt}\footnotesize
\begin{tabular}{lrrrrrr}
\toprule
$\lambda$ & Ent. & $<0.1\%$ & 0.1--0.3\% & 0.3--1\% & mAP & Mass \\
\midrule
$-0.1$  & 0.90 & \textbf{0.908} & \textbf{0.941} & \textbf{0.959} & 0.790 & 0.026 \\
$0$     & 0.44 & 0.899 & 0.933 & 0.951 & \textbf{0.800} & \textbf{0.088} \\
$0.01$  & 0.29 & 0.899 & 0.929 & 0.950 & 0.790 & 0.072 \\
$0.1$   & 0.06 & 0.886 & 0.916 & 0.939 & 0.713 & 0.061 \\
\midrule
mean   & 1.00 & 0.862 & 0.918 & 0.928 & 0.734 & 0.007 \\
oracle & --   & 0.973 & 0.989 & 0.992 & 0.944 & --    \\
\bottomrule
\end{tabular}
\end{table}
\paragraph{Localisation audit.}
Table~\ref{tab:localisation} shows that the head's attention sits on the target far more often
than chance in every bin. For sub-patch targets ($<0.1\%$) the mean attention mass inside the
boxes is 0.088 against a chance level of 0.007, a 12-fold ratio, and the arg-max patch lies
inside a box for 13\% of positives; for targets of 0.1--1\% the ratio rises to 24 and pointing
accuracy to 26--38\%; for large targets the mass inside the box reaches 0.81 and pointing
accuracy 84\%. The absolute mass on sub-patch targets is still small, which is consistent with
Figure~\ref{fig:attention}: on the smallest objects the head often lands one patch beside the
box or on a correlated object (the player rather than the racket), and the box itself covers at
most one of 256 patches. The attention is nevertheless far from uniform: its normalised entropy
settles at 0.43 during training (uniform $=1$), so the head is not re-deriving mean pooling.

\paragraph{Robustness and variants.}
Three seeds of the $\lambda=0$ head give sub-patch AUROC 0.899, 0.891 and 0.889 (mean 0.893,
s.d.\ 0.005) and mAP 0.800, 0.802 and 0.798, so the gain over mean pooling (0.862) is not a
seed effect. Max-MIL with top-8 instead of top-1 pooling recovers on large targets (0.958 vs.\
0.945) and in mAP (0.654 vs.\ 0.631) but stays below mean pooling in every bin: the linear
per-patch scorer is the limitation, not the pooling rule. The multi-scale head
(\texttt{attn\_ms}), which attends jointly over the $16\times16$ grid and its $2\times$ and
$4\times$ average-pooled versions (336 tokens with learned scale embeddings, 478k parameters),
is the best read-out we trained: 0.908 on sub-patch targets, 0.948 and 0.963 in the next two
bins, 0.990 on large targets (level with mean pooling), overall AUROC 0.973 and mAP 0.816, with
attention mass on sub-patch targets 16 times chance and pointing accuracy 0.195 (vs.\ 0.129
for the single-scale head). It closes 41\% of the oracle gap on sub-patch targets, 42\% in the
0.1--0.3\% bin and 55\% in the 0.3--1\% bin, so letting the head choose its scale helps exactly
where the single-scale head was weakest, and it does so while looking at the object more, not
less.

\subsection{Discussion and caveats}
\label{sec:discussion}
\paragraph{Context priors.}
COCO scenes are strongly predictive of their objects: a tiny sports ball co-occurs with a
person swinging a bat, a toothbrush with a sink. A global embedding can therefore score a small
object well above chance without encoding the object at all, simply by encoding the scene
\citep{oliva2007context}. This is why the AUROC curve alone cannot separate ``sees the object''
from ``guesses from context''. Two of our measurements can. The oracle gap tells us how much
more is available when the read-out is pointed at the object; if the global read-outs are
already close to the oracle in a bin, context is doing the work. The localisation audit tells
us whether the attention head is looking at the object or at its companions; a head that beats
mean pooling while placing no more attention mass in the box than chance has learned context,
not the target. On the first count the data say that context does much but not all of the
work: the global read-outs sit 0.11 AUROC below the oracle on sub-patch targets and 0.06 below
on one-to-two-patch targets, so a large share of the available evidence is discarded, while the
0.83--0.86 they retain is far above what a context-free domain would show. On the second count
the head passes the audit: its attention mass inside the box is 12 to 24 times chance on small
targets, so its gain over mean pooling comes with localisation rather than instead of it. The
tennis-racket panel of Figure~\ref{fig:attention} shows that context routes still exist, and
the audit is what makes them visible.

\paragraph{What this does and does not establish.}
On natural images the cliff may be shallower than in chest radiography, where a lesion has no
predictive context and the rest of the image looks the same with or without it. We expect and
say specifically that the drop we observe for global read-outs on COCO, 0.13 to 0.15 AUROC
between the largest and the sub-patch bin, is a graded loss and is not comparable to the
chance-level result in chest radiography \citep{smalllesion2026}; it is a statement about
DINOv2 ViT-B/14 at 224 px on COCO, not about foundation models in general.
What the controlled comparison does establish, if the oracle stays high where the global
read-outs fall, is that for \emph{this} encoder and \emph{these} tokens the loss occurs at
aggregation and is recoverable by a read-out that can select patches.

% ======================================================================
\section{Updated Plan}
\label{sec:plan}
% ======================================================================
Three weeks remain. The plan below keeps the midterm study as the core deliverable and adds the
radiograph replication only after it is complete.

\begin{table}[h]
\centering\small
\caption{Week-by-week plan to the final report. Audio and distillation remain deferred, as
recommended in the pre-proposal feedback; LUNA16 (3D CT) is dropped.}
\label{tab:plan}
\begin{tabular}{p{0.95in}p{2.1in}}
\toprule
Week & Deliverable \\
\midrule
Nov 16--22 &
  Extend the resolution ablation (336 px; the attention head and max-MIL at 448 px) and
  repeat the head over three seeds; add a
  second frozen backbone with a different pretraining objective and patch size (SigLIP or CLIP
  ViT-B/16) on the same harness; per-class analysis separating context-rich from context-poor
  classes; multi-scale head. \\
Nov 23--29 &
  \textbf{Clinic, Nov 23}: present the blindness curve, the oracle gap and the localisation
  audit; get feedback on the causal argument. Start the chest-radiograph replication with the
  same harness: NIH ChestX-ray14 image-level labels for training, ChestX-Det10 boxes for
  binning and the oracle, RAD-DINO and BiomedCLIP as backbones. \\
Nov 30--Dec 6 &
  \textbf{Poster and demo, Nov 30}: live demo of the curve and the attention maps on
  user-chosen images. Then consolidate: repeat key runs over three seeds, finish the
  radiograph curves, write the size-stratified evaluation recommendation. Tiled-teacher
  distillation only if everything above is complete. \\
Dec 7 &
  \textbf{Final report}: cross-dataset (COCO, chest radiographs) and cross-backbone blindness
  curves, the recovery method with its localisation audit, the evaluation recommendation, and
  released code, cached curves and reproduction commands. \\
\bottomrule
\end{tabular}
\end{table}

\paragraph{Risks and fallbacks.}
The midterm result stands on its own, so the remaining risks are to the extensions. The
chest-radiograph replication depends on downloading NIH ChestX-ray14 (42 GB) and on
ChestX-Det10 boxes aligning with it; if either fails we use the 984 boxed images that ship with
ChestX-ray14 for binning only. A second backbone with a 16-pixel patch changes the grid and the
one-patch boundary, which the harness already parameterises. Seeds: every number here is one
seed, so the final report repeats the head and the probes over three seeds before claiming any
difference under 0.01 AUROC. Audio, distillation and further backbones stay out of scope for
this semester unless everything above is complete.

% ======================================================================
\bibliographystyle{unsrtnat}
\bibliography{refs}

\end{document}
